% ---
% id: receiver/rem:overlap
% kind: rem
% title: "Disjoint mode supports as the commutation hypothesis"
% summary: >-
%   The contract's only commutation result, [[prop:mode-algebra]] (iii), needs
%   disjoint mode supports; being unrelated in a declared graph is not enough.
%   Unrelated elements that share a mode can depend on their order, which is
%   why a well-formed graph must order them.
% prominence: standard
% status: accepted
% uses: []
% mentions: [prop:mode-algebra, def:graph, thm:compilation-well-defined]
% assumes: []
% refs: []
% provenance: newly-derived
% lean:
% ---

\begin{remark}\label{rem:overlap}
Proposition~\ref{prop:mode-algebra}\,(iii) is the \emph{only} commutation this
contract has, and it needs disjointness of the \emph{mode supports}, not
disjointness in the wiring diagram. Two elements that are unrelated in a
declared graph may still touch a common mode, and then their product can
depend on the order. That is exactly why Definition~\ref{def:graph} carries
condition~(E3) and why Theorem~\ref{thm:compilation-well-defined} is a theorem rather than a
formality.
\end{remark}
